\documentclass[10pt,twocolumn]{article}

% === Packages ===
\usepackage[letterpaper,margin=0.56in]{geometry}
\usepackage[utf8]{inputenc}
\usepackage[T1]{fontenc}
\usepackage[brazilian]{babel}
\usepackage{amsmath,amssymb,amsthm}
\usepackage{graphicx}
\usepackage{booktabs}
\usepackage{multirow}
\usepackage{natbib}
\usepackage{ragged2e}
\usepackage{threeparttable}
\usepackage{float}
\usepackage{placeins}
\usepackage{lmodern}
\usepackage{caption}
\usepackage{xurl}
\usepackage{hyperref}
\hypersetup{
  colorlinks=true,
  linkcolor=black,
  citecolor=blue,
  urlcolor=black,
  pdftitle={Nota de Politica Publica: Quanto de produtividade precisamos para reduzir a jornada de trabalho?},
  pdfauthor={Victor Rangel}
}
\Urlmuskip=0mu plus 1mu\relax
\graphicspath{{figures/}}
\setcitestyle{authoryear,round}
\captionsetup{font=scriptsize,labelfont=bf}
\setlength{\floatsep}{4pt plus 1pt minus 1pt}
\setlength{\textfloatsep}{5pt plus 1pt minus 1pt}
\setlength{\dblfloatsep}{4pt plus 1pt minus 1pt}
\setlength{\dbltextfloatsep}{5pt plus 1pt minus 1pt}
\setcounter{topnumber}{4}
\setcounter{dbltopnumber}{4}
\renewcommand{\topfraction}{0.95}
\renewcommand{\dbltopfraction}{0.95}
\renewcommand{\textfraction}{0.03}
\renewcommand{\floatpagefraction}{0.85}
\renewcommand{\dblfloatpagefraction}{0.85}
\setlength{\parskip}{1pt}
\setlength{\parindent}{10pt}
\raggedbottom
\emergencystretch=3em
\renewcommand{\abstractname}{Resumo}
\renewcommand{\refname}{Refer\^encias}
\renewcommand{\tablename}{Tabela}
\renewcommand{\figurename}{Figura}
\newcommand{\tabnotes}[1]{\par\vspace{2pt}\noindent\begin{minipage}{\linewidth}\footnotesize #1\end{minipage}}

% Public replication package and online appendix
\newcommand{\onlineappendixurl}{https://github.com/vr-rodrigues/jornada-6x1-replication}

% === Metadata ===
\title{\vspace{-1.35cm}\textbf{Nota de Pol\'itica P\'ublica: Quanto de produtividade precisamos para reduzir a jornada de trabalho?}}
\author{Victor Rangel\textsuperscript{*}\\[-1pt]{\normalsize Insper}}
\date{}

\begin{document}
\twocolumn[
\maketitle
\vspace{-0.55cm}

% === Resumo ===
\begin{abstract}
\noindent
O debate brasileiro sobre redu\c{c}\~ao da jornada deixou de ser apenas uma escolha entre manter 44 horas ou ir direto para 36. Hoje tamb\'em est\~ao na mesa alternativas de 40 horas, escala 5x2 e transi\c{c}\~oes faseadas. Esta nota faz uma pergunta simples para orientar essa escolha: quanto a economia teria de ficar mais produtiva para que cada op\c{c}\~ao n\~ao reduzisse a produ\c{c}\~ao no curto prazo? Para responder, combino dados brasileiros de horas trabalhadas, informalidade, tamanho das empresas e composi\c{c}\~ao setorial com um modelo de ajuste entre emprego formal e informal. O resultado principal \'e que uma redu\c{c}\~ao para 40 horas exige um ganho de produtividade de cerca de $2\%$. A redu\c{c}\~ao direta para 36 horas exige um salto bem maior, entre $6{,}6\%$ e $8{,}2\%$, valor alto frente ao hist\'orico recente da produtividade brasileira. A informalidade tamb\'em sobe no cen\'ario de 36 horas, em torno de $1{,}6$ a $1{,}9$ ponto percentual, mas o principal custo vem de menos horas formais trabalhadas. O exerc\'icio n\~ao diz se a reforma deve ou n\~ao avan\c{c}ar; ele mostra que tamanho, ritmo e instrumentos de transi\c{c}\~ao mudam bastante a conta. Para formuladores de pol\'itica, a mensagem \'e direta: uma rota faseada, com parada perto de 40 horas, exige uma meta de produtividade muito menor do que um salto imediato para 36 horas.
\end{abstract}

\vspace{0.05cm}
\noindent{\footnotesize \textbf{Palavras-chave:} jornada de trabalho, informalidade, produtividade, Brasil, desenho de pol\'itica, transi\c{c}\~ao

\textbf{JEL:} E24, J22, J46, O17, O47.}
\vspace{0.35cm}
]
\begingroup
\renewcommand{\thefootnote}{*}
\footnotetext{E-mail: \href{mailto:victorr@insper.edu.br}{victorr@insper.edu.br}. ORCID: \href{https://orcid.org/0000-0002-4520-2795}{https://orcid.org/0000-0002-4520-2795}.}
\endgroup

\FloatBarrier

% ==================================================================
% 1. Introdu\c{c}\~ao
% ==================================================================
\section{Introdu\c{c}\~ao}\label{sec:intro}

O debate brasileiro sobre jornada hoje envolve mais de um endpoint estatut\'ario. A PEC 8/2025 foi apresentada como uma proposta de quatro dias e 36 horas \citep{CamaraPEC82025Protocol}. Em abril de 2026, a admissibilidade da PEC 8/2025 e da PEC 221/2019 havia sido aprovada na CCJ, uma comiss\~ao especial havia sido criada, e a cobertura da C\^amara reportava uma proposta do relator com alternativas de 40 horas, escala 5x2 e transi\c{c}\~oes faseadas \citep{CamaraFortyHours2025,CamaraAdmissibility2026,CamaraSpecialCommission2026}. A pergunta quantitativa \'e qu\~ao grande teria de ser o ganho de produtividade em cada ponto dessa sequ\^encia de tetos de horas.

Este artigo calcula o ganho \'unico de PTF que manteria o produto agregado inalterado no curto prazo quando a jornada estatut\'aria cai de 44 horas para tetos alternativos entre 44 e 36 horas. Chamamos esse ganho de benchmark de neutralidade em produto e o denotamos $A_{\text{req}}$. O teto de 36 horas \'e o endpoint de alta intensidade da proposta original. O teto de 40 horas tornou-se uma alternativa central no debate p\'ublico. O exerc\'icio n\~ao sup\~oe que a produtividade n\~ao reaja \`a reforma. Ele calcula o tamanho da rea\c{c}\~ao que seria necess\'aria para neutralizar a perda mec\^anica de produto.

O ambiente brasileiro complica o c\'alculo. A informalidade \'e de $37{,}8\%$ pela defini\c{c}\~ao restrita da PNAD e $44{,}2\%$ pela ampla do IBGE, e varia de $41\%$ em servi\c{c}os a $72\%$ na agricultura. Modelos existentes de reformas de jornada tendem a abstrair da informalidade \citep{Cacciatore2016} ou a trat\'a-la por outra margem \citep{Ulyssea2018}. Este artigo aplica essa preocupa\c{c}\~ao a um corte amplo do teto estatut\'ario.

O modelo \'e um arcabou\c{c}o est\'atico de equil\'ibrio parcial calibrado com microdados brasileiros. Ele combina um agregador CES sobre trabalho efetivo formal e informal, uma fun\c{c}\~ao $e(h)$ de produtividade por hora com pico em $40$ horas, dois grupos por tamanho de firma e prefer\^encias GHH para bem-estar. A elasticidade de substitui\c{c}\~ao \'e disciplinada por uma janela de pr\^emio salarial formal-informal brasileiro de \citet{MeghirNaritaRobin2015}. Essa disciplina n\~ao identifica um ponto \'unico: os resultados principais reportam tamb\'em o envelope conjunto de $(\sigma,\omega,\eta_I)$. O exerc\'icio deve ser lido como contabilidade macro disciplinada por momentos, n\~ao como um modelo completo de equil\'ibrio geral da reforma.

O objeto central \'e $A_{\text{req}}$ como fun\c{c}\~ao do teto. Uma redu\c{c}\~ao de 44 para 40 horas exige cerca de $2\%$ de PTF e deixa positivo o indicador GHH representativo. O endpoint de 36 horas \'e mais exigente. Na especifica\c{c}\~ao de \emph{fadiga s\'o acima de 40h}, o endpoint tem $A_{\text{req}}=6{,}63\%$. Na especifica\c{c}\~ao de \emph{penalidade bilateral}, em que desvios para cima e para baixo de 40h reduzem a efici\^encia hor\'aria, $A_{\text{req}}=8{,}18\%$. O envelope conjunto $(\sigma,\omega,\eta_I)$ \'e $[5{,}62\%,\,7{,}48\%]$ e $[6{,}93\%,\,9{,}23\%]$, respectivamente. A informalidade agregada sobe $+1{,}6$ a $1{,}9$ p.p. no endpoint de 36 horas. O canal principal do custo de produto \'e a redu\c{c}\~ao mec\^anica de horas, n\~ao a realoca\c{c}\~ao formal-informal. Como benchmark de escala, a PWT 11.0 mostra que a PTF brasileira caiu no horizonte p\'os-1990; portanto, o endpoint de 36 horas exige um ganho de produtividade grande em compara\c{c}\~ao com a experi\^encia recente brasileira.

Em rela\c{c}\~ao \`a literatura, o exerc\'icio traz a margem formal-informal para uma an\'alise estrutural de reforma de jornada. Essa margem \'e modesta para o produto agregado sob par\^ametros disciplinados pelo pr\^emio salarial, mas muda a incid\^encia da reforma. O benchmark $A_{\text{req}}$ tamb\'em permite comparar cada teto de horas com o hist\'orico brasileiro de PTF. O intervalo para $\sigma_{\text{sub}}$ deve ser lido como disciplina de conjunto, n\~ao como estimativa pontual da elasticidade relevante para qualquer texto legal espec\'ifico.

Este artigo n\~ao estima o efeito causal da PEC 8/2025, da PEC 221/2019 ou de qualquer texto espec\'ifico que venha a sair da comiss\~ao especial. N\~ao modela equil\'ibrio geral, barganha salarial, sele\c{c}\~ao de trabalhadores, ajuste de pre\c{c}os ou aprofundamento de capital como canais centrais. O exerc\'icio calcula o ganho de PTF que tornaria tetos alternativos de horas neutros em produto sob par\^ametros disciplinados por dados brasileiros. Evid\^encia de Portugal 1996 \citep{AsaiLopesTondini2024} e Brasil 1988 \citep{GonzagaMenezesFilhoCamargo2003} entra como checagem direcional de plausibilidade, n\~ao como limite inferior identificado para o $A_{\text{req}}$ efetivo brasileiro.

Assim, o artigo responde a uma pergunta de escala: qual ganho de produtividade tornaria cada teto de horas neutro em produto em um ambiente de curto prazo com informalidade e emprego fixo? Ele n\~ao responde se a reforma geraria esse ganho, como sal\'arios e pre\c{c}os se ajustariam em equil\'ibrio geral, nem qual texto legal maximizaria bem-estar social. Essa distin\c{c}\~ao orienta a leitura dos resultados de produto, informalidade e bem-estar abaixo.

% ==================================================================
% 2. Fatos brasileiros e arcabou\c{c}o
% ==================================================================
\section{Fatos brasileiros e modelo}\label{sec:framework}

As horas formais est\~ao concentradas no teto estatut\'ario. O baseline mapeia a distribui\c{c}\~ao de horas contratadas do DIEESE em tr\^es bins do modelo. Os $8{,}5\%$ dos trabalhadores formais na categoria $\leq 39$h entram no bin de 36h, os $26{,}9\%$ em 40h permanecem no bin de 40h, e os $64{,}6\%$ na categoria $\geq 41$h entram no bin de 44h \citep{DIEESE2024}. A PNAD Cont\'inua Q4 2024, baseada em horas habituais, mostra uma distribui\c{c}\~ao mais achatada, com m\'edia de $41{,}2$ \citep{PNAD2025}. A maioria dos trabalhadores formais opera em 40 horas ou acima, que \'e a margem restringida por um teto de 36h.

A informalidade \'e grande e dispersa. A taxa nacional \'e de $37{,}8\%$ pela defini\c{c}\~ao restrita da PNAD e $44{,}2\%$ pela ampla do IBGE \citep{PNAD2025}. A faixa setorial vai de $41{,}2\%$ em servi\c{c}os a $71{,}8\%$ na agricultura (Tabela~\ref{tab:facts_informality}). O pr\^emio salarial formal-informal tamb\'em disciplina o problema. \citet{MeghirNaritaRobin2015} documentam um pr\^emio salarial condicional $R\in[1{,}15,\,1{,}55]$, medido como raz\~ao entre o sal\'ario-hora m\'edio formal e o sal\'ario-hora m\'edio informal entre trabalhadores compar\'aveis. A calibra\c{c}\~ao usa esse intervalo para ancorar a elasticidade CES de substitui\c{c}\~ao $\sigma_{\text{sub}}$ entre trabalho formal e informal.

A PTF brasileira permaneceu estagnada ou em queda ao longo do per\'iodo p\'os-1990. A s\'erie \texttt{rtfpna} da Penn World Table 11.0 mostra queda anualizada de $-0{,}83\%$ ao ano de 1990 a 2023 e $-0{,}82\%$ no horizonte pr\'e-Covid de 1990 a 2019 \citep{FeenstraInklaarTimmer2015,PWT2025,FREDPWTBrazil2026}. A melhor d\'ecada m\'ovel p\'os-1990 na s\'erie atualizada \'e de 2000 a 2010, com apenas $+0{,}05\%$ ao ano. O benchmark de PTF na Se\c{c}\~ao~\ref{sec:results} usa essa s\'erie atualizada como \^ancora para $A_{\text{req}}$.

Esses fatos pedem um modelo com tr\^es margens. A primeira \'e a concentra\c{c}\~ao de trabalhadores formais perto do teto estatut\'ario. A segunda \'e a substitui\c{c}\~ao entre trabalho efetivo formal e informal, disciplinada pelo pr\^emio salarial observado. A terceira \'e a produtividade por hora, que varia com a jornada. O modelo combina essas margens com dois grupos por tamanho de firma para incid\^encia diferencial e prefer\^encias GHH para avalia\c{c}\~ao de bem-estar.

\begin{table*}[t]
\centering
\caption{Momentos setoriais e nacionais do mercado de trabalho, PNAD Cont\'inua Q4 2024.}\label{tab:facts_informality}
\small
\begin{tabular}{lcccccc}
\hline\hline
Setor & $\lambda_s$ & Informalidade & Hrs m\'edias & $\theta_{36}$ & $\theta_{40}$ & $\theta_{44}$ \\
 & & (ampla) & (formal) & & & \\
\hline
Agricultura & 0{,}076 & 71{,}8\% & 44{,}1 & 0{,}101 & 0{,}269 & 0{,}630 \\
Ind\'ustria & 0{,}205 & 44{,}4\% & 42{,}5 & 0{,}048 & 0{,}417 & 0{,}535 \\
Servi\c{c}os & 0{,}720 & 41{,}2\% & 40{,}8 & 0{,}171 & 0{,}410 & 0{,}419 \\
\hline
Nacional    & 1{,}000 & 44{,}2\% & 41{,}2 & 0{,}143 & 0{,}406 & 0{,}451 \\
\hline\hline
\end{tabular}
\\[2pt]
\begin{minipage}{0.95\textwidth}
\justifying
{\footnotesize \textit{Notas:} $\lambda_s$ \'e a participa\c{c}\~ao do setor no emprego nacional. Informalidade usa a defini\c{c}\~ao ampla do IBGE (VD4009 combinada com o indicador V4017 CNPJ). ``Hrs m\'edias'' \'e a m\'edia de horas semanais dos trabalhadores \emph{formais} (V4039, habitual). $\theta_{36}$, $\theta_{40}$, $\theta_{44}$ s\~ao as parcelas de trabalhadores formais em cada bin de horas contratadas. \textit{Fonte:} PNAD Cont\'inua Q4 2024, microdados IBGE, ponderados pela amostra, $N=209{,}311$ \citep{PNAD2025}.}
\end{minipage}
\end{table*}

Uma firma representativa no grupo $g\in\{\text{Pequena},\text{Grande}\}$ produz o produto $Y_g$ a partir do capital $K_g$ e do trabalho efetivo $L_g$ com tecnologia Cobb-Douglas. Firmas pequenas t\^em at\'e 49 empregados. Firmas grandes t\^em 50 ou mais.
\begin{equation}\label{eq:production}
Y_g = A\cdot K_g^{\alpha}\cdot L_g^{1-\alpha},
\end{equation}
em que $A$ \'e a produtividade agregada e $\alpha$ \'e a parcela do capital na renda. O trabalho efetivo $L_g$ combina trabalhadores formais ($F$) e informais ($I$) por um agregador CES:
\begin{equation}\label{eq:ces}
L_g = \left[\omega L_{F,g}^{\rho} + (1-\omega)L_{I,g}^{\rho}\right]^{1/\rho},\quad \rho=(\sigma_{\text{sub}}-1)/\sigma_{\text{sub}},
\end{equation}
em que $\omega$ \'e o peso CES do trabalho formal (um $\omega$ maior torna o trabalho formal mais importante na produ\c{c}\~ao) e $\sigma_{\text{sub}}$ \'e a elasticidade de substitui\c{c}\~ao entre trabalho formal e informal. Quando $\sigma_{\text{sub}}<1$, os dois insumos s\~ao complementos, de modo que um aumento do custo do trabalho formal \emph{n\~ao} provoca grande realoca\c{c}\~ao para a margem informal. Os dois insumos de trabalho s\~ao
\[
L_{F,g}=N_{F,g}\sum_b\theta_b h_b e(h_b),\qquad L_{I,g}=\eta_I N_{I,g} h_I e(h_I),
\]
em que $N_{F,g}$ e $N_{I,g}$ s\~ao os quantitativos de trabalhadores formais e informais no grupo $g$ (com $N_{I,g}\equiv N_g-N_{F,g}$ e o emprego total do grupo $N_g$ fixo), $h_b$ indexa os tr\^es bins de horas contratadas $b\in\{36,40,44\}$ com parcelas $\theta_b$ entre os formais, $h_I=44$ s\~ao as horas semanais informais sem teto estatut\'ario, e $\eta_I$ \'e a produtividade por cabe\c{c}a de um trabalhador informal relativa \`a de um formal. A efici\^encia por hora segue $e(h)=\exp\{-\kappa(h-h^*)^2\}$ acima do pico $h^*=40$, em que $\kappa$ \'e a curvatura da penalidade de fadiga que jornadas longas imp\~oem \`a produtividade. Abaixo do pico, o artigo considera duas hip\'oteses transparentes. A especifica\c{c}\~ao preferida, chamada \emph{fadiga s\'o acima de 40h}, fixa $e(h)=1$ para $h\leq h^*$: jornadas longas cansam, mas jornadas mais curtas n\~ao recebem penalidade mec\^anica. A alternativa conservadora, chamada \emph{penalidade bilateral}, mant\'em a mesma forma quadr\'atica abaixo de $h^*$: desvios para cima ou para baixo de 40h reduzem a efici\^encia por hora. A primeira hip\'otese \'e informada por \citet{Pencavel2015}, que documenta a perna de fadiga acima do pico mas n\~ao caracteriza a efici\^encia abaixo dele. A segunda serve como limite superior conservador para o ganho de PTF requerido.

Cada firma escolhe o quantitativo de trabalhadores formais $N_{F,g}$ para maximizar o lucro l\'iquido de tr\^es cunhas de custo:
\begin{equation}\label{eq:firm}
Y_g-\tau_g N_{F,g}-\tfrac{1}{2}\pi_m N_{I,g}^2-\tfrac{1}{2}\gamma_{F,g}(N_{F,g}-\bar{N}_{F,g})^2,
\end{equation}
em que $\tau_g$ \'e um custo linear por trabalhador formal (obriga\c{c}\~oes CLT/FGTS e enforcement dependente do tamanho), $\pi_m$ \'e um custo convexo de informalidade (multas, risco reputacional e acesso limitado a insumos do setor formal) e $\gamma_{F,g}$ \'e um custo quadr\'atico de ajuste que penaliza desvios em rela\c{c}\~ao ao quantitativo formal pr\'e-reforma $\bar{N}_{F,g}$. A cunha $\tau_g$ \'e calibrada para casar a informalidade do grupo e desempenha o papel de uma distor\c{c}\~ao dependente do tamanho no esp\'irito de \citet{HsiehKlenow2009} e \citet{RestucciaRogerson2008}, com microfunda\c{c}\~oes de formalidade end\'ogena em \citet{DixCarneiroGoldbergMeghirUlyssea2026}. Cortar o teto estatut\'ario de $h_0=44$ horas para $h_1$ reduz o trabalho efetivo tanto via os bins de horas $h_b$ quanto via a fun\c{c}\~ao de efici\^encia $e(h_b)$, e alarga a cunha entre o produto marginal do trabalho formal e seu custo privado.

O objeto de interesse \'e o multiplicador de PTF que restaura o produto agregado:
\begin{equation}\label{eq:areq}
A_{\text{req}}:\quad \sum_g Y_g\!\left(A\,(1+A_{\text{req}}),K_g,h_1\right)=\sum_g Y_g(A,K_g,h_0).
\end{equation}
Em palavras, $A_{\text{req}}$ \'e o ganho percentual \'unico na produtividade agregada $A$ que compensaria exatamente a perda de produto de curto prazo induzida pela mudan\c{c}a do teto de $h_0$ para $h_1$. Trata-se de um \emph{benchmark} cont\'abil que o leitor pode comparar com o hist\'orico de PTF brasileiro, n\~ao de uma taxa de crescimento prevista. Para bem-estar, as prefer\^encias s\~ao GHH \citep{GreenwoodHercowitzHuffman1988} com $\sigma_{\text{ghh}}=1$ e elasticidade de Frisch $1/\nu=0{,}5$. A varia\c{c}\~ao compensat\'oria \'e
\begin{equation}\label{eq:cv}
\Delta\text{CV}=\frac{C_1-\psi h_1^{1+\nu}/(1+\nu)}{C_0-\psi h_0^{1+\nu}/(1+\nu)}-1.
\end{equation}
O numerador e o denominador s\~ao os equivalentes de consumo de bem-estar p\'os- e pr\'e-reforma sob prefer\^encias GHH. Eles medem consumo l\'iquido da desutilidade do trabalho monetizada $\psi h^{1+\nu}/(1+\nu)$. $\Delta$CV \'e a varia\c{c}\~ao percentual de bem-estar. Um valor positivo indica que a reforma melhora o bem-estar. Um valor negativo indica que o reduz.

O modelo \'e est\'atico e de equil\'ibrio parcial. O capital \'e pr\'e-determinado, o emprego total $N$ \'e fixo, e as retroalimenta\c{c}\~oes de equil\'ibrio geral (sal\'ario-pre\c{c}o, monet\'aria, fiscal, intersetorial) est\~ao ausentes. O ap\^endice online enumera os canais modelados, bracketeados e omitidos, incluindo compliance, horas extras, barganha salarial, oferta de trabalho e liga\c{c}\~oes setoriais. Uma elasticidade de custo de uso $\varepsilon_K\in[0{,}2,\,0{,}4]$ \citep{Caballero1999} com parcela de capital $\alpha=0{,}35$ implica
\begin{equation}\label{eq:pe_bound}
\Delta A_{\text{req}}^{K}\approx -\alpha\cdot\varepsilon_K\cdot|\Delta L_F/L|\in[-1{,}0,\,-1{,}8]\text{ p.p.},
\end{equation}
onde $\Delta A_{\text{req}}^{K}$ \'e a corre\c{c}\~ao para baixo do $A_{\text{req}}$ de curto prazo uma vez incorporada a resposta end\'ogena do estoque de capital. O an\'alogo de m\'edio prazo do resultado principal cai $1$ a $2$ p.p.\ em horizonte de $3$ a $5$ anos.

A Tabela~\ref{tab:params} lista os par\^ametros. Par\^ametros externos v\^em da PWT, DIEESE, PNAD Cont\'inua e RAIS \citep{PNAD2025,RAIS2022,DIEESE2024,PWT2025}. Par\^ametros calibrados por momentos ($\tau_g$, $\pi_m$, $\kappa$, $\psi$) s\~ao ajustados internamente. Par\^ametros escolhidos pelo autor ($\eta_I$, $e_q$, $h^*$, $h_I$, $\gamma_F^g$) situam-se em faixas consistentes com a evid\^encia sobre horas de trabalho em \citet{Pencavel2015} e \citet{CollewetSauermann2017}, a literatura de produtividade informal \citep{LaPortaShleifer2014} e regras de compliance de folha \citep{BrasilLei8212,ReceitaESocialRural2021}. A evid\^encia transnacional sobre horas em \citet{BickFuchsSchundelnLagakos2018} fornece suporte auxiliar para produtividade marginal decrescente. O ap\^endice online apresenta fontes e sensibilidade completas para cada par\^ametro.

\begin{table*}[t]
\centering
\caption{Valores de par\^ametros por base evidencial.}\label{tab:params}
\small
\begin{tabular}{@{}c p{3.1cm} c p{6.1cm}@{}}
\hline\hline
S\'imbolo & Descri\c{c}\~ao & Valor & Fonte \\
\hline
\multicolumn{4}{l}{\textit{Externos}} \\
\hline
$\alpha$     & Parcela de capital             & 0{,}35   & PWT + \citet{Gollin2002} \\
$\omega$     & Peso formal CES                & 0{,}622  & $1-$ informalidade restrita \citep{PNAD2025} \\
$\theta_b$   & Parcelas por bin de horas      & $(0{,}085, 0{,}269, 0{,}646)$ & Mapeadas de categorias de horas contratadas do DIEESE \citep{DIEESE2024} \\
$\nu$        & Inverso da elasticidade de Frisch & 2{,}0 & Padr\~ao (Frisch $=0{,}5$) \\
\hline
\multicolumn{4}{l}{\textit{Calibrados por momentos}} \\
\hline
$\sigma_{\text{sub}}$ & Elasticidade de substitui\c{c}\~ao CES & \textbf{1{,}326} $[1{,}116,\,1{,}469]$ & Disciplina por pr\^emio salarial \citep{MeghirNaritaRobin2015}; ponte PNAD diagn\'ostica \citep{PNAD2025} \\
$\tau_S,\tau_L$       & Cunhas formais (pequena, grande) & 2{,}609;\,0{,}004 & Bisse\c{c}\~ao para informalidade do grupo \\
$\pi_m^S,\pi_m^L$     & Custo convexo de informalidade   & 0{,}000;\,42{,}11 & Bisse\c{c}\~ao \\
$\kappa$               & Curvatura de efici\^encia & $2{,}11\times 10^{-3}$ & Derivado de $e_q$ em $h_{\text{ref}}$ \\
$\psi$                 & Peso de desutilidade GHH & $5{,}74\times 10^{-5}$ & CPO de base \\
\hline
\multicolumn{4}{l}{\textit{Escolhidos pelo autor (dentro de faixas consistentes com a evid\^encia)}} \\
\hline
$\eta_I$      & Efici\^encia relativa informal & 0{,}40   & Ponto m\'edio da raz\~ao salarial LPS (informal/formal) aprox.\ 1/3 a 1/2 \\
$e_q$          & Elasticidade horas--produto    & 0{,}60   & Faixa \citet{Pencavel2015} \\
$h^*$          & Horas de efici\^encia de pico  & 40       & Informado por evid\^encia de horas longas em \citet{Pencavel2015} \\
$h_I$          & Horas informais                & 44       & Sem teto estatut\'ario no setor informal \\
$\gamma_F^S, \gamma_F^L$ & Custos de ajuste         & 0{,}12;\,0{,}03 & RAIS + regras de folha \citep{RAIS2022,BrasilLei8212,ReceitaESocialRural2021} \\
\hline\hline
\end{tabular}
\end{table*}

A elasticidade $\sigma_{\text{sub}}$ \'e disciplinada pelo pr\^emio salarial, n\~ao estimada diretamente. Da CPO do CES, o pr\^emio salarial formal-informal $R$ \'e fun\c{c}\~ao mon\'otona de $\sigma_{\text{sub}}$. Substituindo uma janela ampliada em torno das estimativas de \citeauthor{MeghirNaritaRobin2015}, $R\in[1{,}15,\,1{,}55]$, na CPO do CES no recalibrado $(\omega,\eta_I)=(0{,}622,\,0{,}40)$ obt\'em-se o intervalo
\begin{equation}\label{eq:sigma_interval}
\begin{aligned}
R\in[1{,}15,\,1{,}55]&\Rightarrow
\sigma_{\text{sub}}\in[1{,}116,\,1{,}469],\\
A_{\text{req}}&\in[7{,}50\%,\,8{,}58\%]
\quad\text{(bilateral).}
\end{aligned}
\end{equation}
O valor central $\sigma=1{,}326$ corresponde a $R\approx 1{,}40$.\footnote{Adotamos $R\in[1{,}15,\,1{,}55]$ como janela plaus\'ivel em torno do pr\^emio salarial formal-informal condicional em \citet{MeghirNaritaRobin2015}. O $R$ de MNR \'e um pr\^emio estrutural estimado dentro do modelo de busca-e-pareamento dos autores. Invert\^e-lo via a CPO do CES sup\~oe que o sal\'ario relativo identificado por MNR coincide com a produtividade marginal relativa de trabalho efetivo formal vs.\ informal no agregador CES. O ap\^endice online reporta uma ponte descritiva PNAD/SIDRA com rendimentos mensais publicados por par formal/informal. O par de empregados privados fornece $R=1{,}24$, enquanto o par de trabalhadores dom\'esticos fornece $R=1{,}58$. Esses n\'umeros n\~ao s\~ao regress\~oes com controles, mas mostram que os gaps brasileiros publicados est\~ao pr\'oximos da janela estrutural MNR. O mapeamento de $R$ para $\sigma$ depende de $(\omega,\eta_I)$. As \^ancoras $\omega=0{,}622$ ($= 1$ menos a informalidade restrita PNAD 2025) e $\eta_I=0{,}40$ (ponto m\'edio da raz\~ao salarial informal/formal LPS, $\approx 1/3$ a $1/2$) d\~ao o mapeamento acima. Sob concorr\^encia perfeita e salarista, a identifica\c{c}\~ao da produtividade marginal no CES usa a raz\~ao salarial, n\~ao a raz\~ao de VA por trabalhador, que \'e contaminada por gaps de capital por trabalhador.} O valor central \'e um ponto de refer\^encia expositivo, n\~ao uma elasticidade identificada. Por isso, as principais tabelas reportam tamb\'em o envelope da caixa conjunta $(\sigma,\omega,\eta_I)$. Permitindo que $\omega$ e $\eta_I$ variem em suas faixas consistentes com os dados ($\omega\in[0{,}58,\,0{,}66]$, $\eta_I\in[0{,}33,\,0{,}50]$), o envelope sob penalidade bilateral \'e $[6{,}93\%,\,9{,}23\%]$ nos oito cantos. Sob fadiga s\'o acima de 40h, o $A_{\text{req}}$ central \'e $6{,}63\%$ e o envelope \'e $[5{,}62\%,\,7{,}48\%]$. O ap\^endice online reporta as grades completas. A pequena elasticidade de realoca\c{c}\~ao documentada por \citet{DerenoncourtGerardLagosMontialoux2025} sustenta um $\sigma$ pequeno do tipo intrafirma, distinto das elasticidades altas entre tipos de trabalhadores em estudos de refugiados. O ap\^endice online discute a distin\c{c}\~ao.

A calibra\c{c}\~ao do $\sigma_{\text{sub}}$ disciplinada pelo pr\^emio salarial substitui valores escolhidos pelo autor em rodadas preliminares do artigo. A mudan\c{c}a afeta principalmente a incid\^encia e o sinal da resposta de informalidade. Sob $\sigma_{\text{sub}}<1$, trabalho formal e informal s\~ao complementos e um teto formal mais baixo aproxima os insumos em vez de afast\'a-los. Sob $\sigma_{\text{sub}}>1$, como disciplinado aqui pela janela MNR, eles s\~ao substitutos imperfeitos e o teto empurra trabalhadores marginais formais para a informalidade. A resposta agregada de informalidade no contrafactual sobe. Para o produto agregado, por\'em, a margem dominante continua sendo a redu\c{c}\~ao mec\^anica de horas formais: variar $\sigma$, $\omega$ e $\eta_I$ dentro da caixa disciplinada desloca o endpoint em cerca de dois pontos percentuais, enquanto o corte bruto de horas explica a maior parte da perda.

Reconstruir a calibra\c{c}\~ao com horas habituais da PNAD ($\theta=(0{,}143,\,0{,}406,\,0{,}451)$, informalidade $44{,}2\%$) resulta em $A_{\text{req}}=7{,}80\%$ sob a penalidade bilateral e aproximadamente $5{,}9\%$ sob a fadiga s\'o acima de 40h. O gap \'e maior na especifica\c{c}\~ao de fadiga s\'o acima de 40h porque a distribui\c{c}\~ao $\theta$ da PNAD coloca mais massa formal abaixo do pico de efici\^encia $h^*=40$, que essa especifica\c{c}\~ao trata como neutra enquanto a penalidade bilateral penaliza. Adotamos horas contratadas do DIEESE como baseline porque o teto estatut\'ario \'e um conceito de horas contratadas.

% ==================================================================
% 3. Valida\c{c}\~ao
% ==================================================================
\section{Checagens da calibra\c{c}\~ao e benchmarks direcionais}\label{sec:validation}

O modelo casa os targets de calibra\c{c}\~ao por constru\c{c}\~ao. Os benchmarks externos abaixo s\~ao checagens direcionais de plausibilidade, n\~ao valida\c{c}\~ao causal do contrafactual. As magnitudes da reforma portuguesa de 1996 e da reforma brasileira de 1988 devem ser comparadas em sinal e ordem, n\~ao em n\'ivel.

A Tabela~\ref{tab:calibration_fit} re\'une os targets, checagens dom\'esticas fora da calibra\c{c}\~ao principal e benchmarks externos. Os targets do Painel A s\~ao casados por constru\c{c}\~ao. O Painel B1 compara momentos observados no Brasil com momentos implicados pelo modelo. O modelo fica pr\'oximo nas horas formais m\'edias e na parcela acima de $36$h. Ele subestima a amplitude setorial das horas formais, uma limita\c{c}\~ao da distribui\c{c}\~ao nacional parcimoniosa de horas.

Os momentos do Painel B1 n\~ao s\~ao consequ\^encias alg\'ebricas da calibra\c{c}\~ao. O pacote de replica\c{c}\~ao perturba a distribui\c{c}\~ao formal de horas $\theta$ em uma faixa de $\pm 5$ p.p.\ em torno dos pesos do DIEESE 2024 e re-roda a calibra\c{c}\~ao. Os momentos ajustados do Painel B1 deslocam-se em $0{,}4$ a $0{,}6$ h/semana e $1$ a $2$ p.p. O sinal do gap permanece, mas a magnitude responde \`a perturba\c{c}\~ao. Esses momentos fornecem uma checagem informativa fora das restri\c{c}\~oes de calibra\c{c}\~ao.

O Painel B2 compara a reforma portuguesa de 1996 no teto $44\to 40$h e a reforma brasileira de 1988 no teto $48\to 44$h. Os sinais casam. O produto agregado cai e a produtividade hor\'aria sobe. As magnitudes n\~ao casam em n\'ivel. A estimativa portuguesa a n\'ivel de firma de $+4{,}4\%$ na produtividade hor\'aria est\'a cerca de $2{,}6$ pontos acima da previs\~ao do modelo agregado, uma diferen\c{c}a consistente com sele\c{c}\~ao de trabalhadores, intensifica\c{c}\~ao e aprofundamento de capital setorial que o modelo n\~ao codifica. A evid\^encia brasileira de 1988 aponta na mesma dire\c{c}\~ao, com resposta emp\'irica hor\'aria de sal\'arios ligeiramente maior. O Painel B2 \'e um benchmark direcional de plausibilidade, n\~ao um limite inferior brasileiro para o $A_{\text{req}}$ efetivo. Portugal e o Brasil de 1988 diferem do ambiente da PEC 8/2025 em informalidade, arranjo institucional e magnitude da reforma.

\begin{table*}[t]
\centering
\caption{Checagens da calibra\c{c}\~ao e benchmarks direcionais.}\label{tab:calibration_fit}
\small
\begin{tabular}{lccc}
\hline\hline
Momento & Observado & Modelo & Gap \\
\hline
\multicolumn{4}{l}{\textit{Painel A: Targets de calibra\c{c}\~ao (casados por constru\c{c}\~ao)}} \\
\hline
Informalidade, pequenas firmas ($\leq 49$ empl.)  & 0{,}500 & 0{,}500 & 0{,}000 \\
Informalidade, grandes firmas                     & 0{,}200 & 0{,}200 & 0{,}000 \\
Parcela de emprego formal, pequenas firmas        & 0{,}590 & 0{,}590 & 0{,}000 \\
Parcela de emprego formal, grandes firmas         & 0{,}410 & 0{,}410 & 0{,}000 \\
Informalidade agregada, PNAD restrita             & 0{,}378 & 0{,}377 & 0{,}001 \\
Pr\^emio salarial $R$, MNR (2015) target central  & 1{,}400 & 1{,}400 & 0{,}000 \\
\hline
\multicolumn{4}{l}{\textit{Painel B1: Checagens dom\'esticas fora da calibra\c{c}\~ao}} \\
\hline
Horas formais semanais m\'edias, medida habitual           & 41{,}23 & 42{,}24 & 1{,}01 h/sem \\
Parcela formal acima de $36$h, medida habitual (\%)        & 85{,}7  & 91{,}5  & 5{,}80 p.p. \\
Amplitude setorial das horas formais m\'edias              & 3{,}31  & 0{,}80  & $-2{,}51$ h/sem \\
\hline
\multicolumn{4}{l}{\textit{Painel B2: Benchmarks direcionais externos}} \\
\hline
$\Delta Y$ no teto $44\to 40$h, Portugal 1996 (firma)     & $-3{,}20$ & $-2{,}00$ & 1{,}20 p.p. \\
$\Delta Y/h$ no teto $44\to 40$h, Portugal 1996 (\%)      & $+4{,}40$ & $+1{,}77$ & 2{,}63 p.p. \\
$\Delta W/h$ em $48\to 44$h, Brasil 1988 (\%, GMC 2003)   & $+7{,}00$ & $+1{,}77$ & 5{,}23 p.p. \\
\hline\hline
\end{tabular}
\\[2pt]
\begin{minipage}{0.95\textwidth}
\justifying
{\footnotesize Painel B1 linhas 1 e 2 comparam momentos de horas habituais da PNAD com momentos do modelo implicados pela calibra\c{c}\~ao de horas contratadas. Linha 3 reporta a amplitude entre agricultura, ind\'ustria e servi\c{c}os nas horas formais semanais m\'edias. O valor do modelo vem da extens\~ao setorial com a mesma estrutura de bins de horas, de modo que o gap negativo significa que o modelo subestima a dispers\~ao setorial observada nas horas formais. O gap de horas agregadas \'e movido principalmente pela hip\'otese mantida $h_I=44$ para trabalhadores informais e \'e tratado como margem de sensibilidade no ap\^endice online, n\~ao como alvo de valida\c{c}\~ao ($A_{\text{req}}$ desloca-se em menos de $0{,}2$ p.p.\ para $h_I\in[38,\,46]$). Painel B2 compara a previs\~ao agregada do modelo no teto $44\to 40$h com a estimativa portuguesa a n\'ivel de firma \citep{AsaiLopesTondini2024} e com a reforma brasileira de 1988 de magnitude compar\'avel (48$\to$44h, $\sim$8\%) documentada por \citet{GonzagaMenezesFilhoCamargo2003}. Ambos os benchmarks externos sugerem que o modelo \'e conservador na resposta de produtividade hor\'aria, como esperado porque o $e(h)$ captura apenas fadiga, n\~ao sele\c{c}\~ao de trabalhadores, intensifica\c{c}\~ao ou aprofundamento setorial de capital. Fontes: PNAD/SIDRA, RAIS 2022, DIEESE 2024, \citet{MeghirNaritaRobin2015}, \citet{AsaiLopesTondini2024} e \citet{GonzagaMenezesFilhoCamargo2003} \citep{PNAD2025,RAIS2022,DIEESE2024}.}
\end{minipage}
\end{table*}

% ==================================================================
% 4. Resultados principais
% ==================================================================
\section{A aritm\'etica de produtividade de tetos alternativos de jornada}\label{sec:results}

O teto de 40 horas e o endpoint de 36 horas s\~ao objetos econ\^omicos diferentes. Mover de 44 para 40 horas corta a cauda longa da distribui\c{c}\~ao contratada e aproxima muitos trabalhadores do pico de efici\^encia assumido. Mover abaixo de 40 horas j\'a n\~ao \'e sobretudo aparar fadiga; \'e retirar horas que o modelo trata como efetivas. Por isso a curva da Figura~\ref{fig:areq_curve} dobra para cima em vez de cair suavemente.

A Tabela~\ref{tab:main_results} mostra o endpoint de 36 horas, mas a compara\c{c}\~ao importante \'e a inclina\c{c}\~ao do caminho. A alternativa de 40 horas exige apenas cerca de $2\%$ de PTF para restaurar o produto. O endpoint de 36 horas exige uma resposta maior: $6{,}63\%$ sob a hip\'otese de fadiga s\'o acima de 40h e $8{,}18\%$ sob a penalidade bilateral. Essa diferen\c{c}a n\~ao \'e mero detalhe t\'ecnico. Na primeira coluna, jornadas mais curtas n\~ao recebem penalidade direta de efici\^encia depois que as horas caem abaixo de 40. Na segunda, afastar-se de 40 horas em qualquer dire\c{c}\~ao tem custo de efici\^encia.

A margem formal-informal se move na dire\c{c}\~ao esperada, mas n\~ao \'e a vil\~a do resultado agregado. A economia calibrada desloca alguns empregos marginais para a informalidade, e isso importa para incid\^encia. Ainda assim, a maior parte da perda de produto vem da aritm\'etica simples de menos horas formais. O envelope conjunto $(\sigma,\omega,\eta_I)$ refor\c{c}a esse ponto: mudar substitui\c{c}\~ao formal-informal, peso formal no CES e efici\^encia informal altera o requisito no endpoint, mas n\~ao transforma 36 horas em um problema pequeno de produtividade. O ajuste de capital pode reduzir cerca de um a dois pontos percentuais do benchmark de curto prazo em horizonte de tr\^es a cinco anos, mas n\~ao elimina a diferen\c{c}a entre 40 e 36 horas.

\begin{table*}[t]
\centering
\caption{Resultados principais no endpoint de 36 horas.}\label{tab:main_results}
\small
\begin{tabular}{@{}l c c@{}}
\hline\hline
 & Fadiga s\'o & Penalidade \\
 & acima de 40h & bilateral \\
\hline
\multicolumn{3}{l}{\textit{Previs\~oes positivas (sem compensa\c{c}\~ao de PTF)}} \\
Produto, $\Delta Y$                       & $-6{,}60\%$       & $-8{,}02\%$ \\
Bem-estar (GHH), $\Delta$CV               & $-1{,}76\%$       & $-3{,}63\%$ \\
Informalidade, $\Delta$Inf                & $+1{,}57$ p.p.    & $+1{,}92$ p.p. \\
Produto por hora, $\Delta Y/h$            & $+2{,}39\%$       & $+0{,}75\%$ \\
\hline
\multicolumn{3}{l}{\textit{Benchmark: ganho de PTF que restaura o produto}} \\
$A_{\text{req}}$, central                 & $6{,}63\%$        & $8{,}18\%$ \\
$A_{\text{req}}$, envelope conjunto       & $[5{,}62;\,7{,}48]\%$ & $[6{,}93;\,9{,}23]\%$ \\
\hline\hline
\end{tabular}
\\[2pt]
\begin{minipage}{0.95\textwidth}
\justifying
{\footnotesize O envelope conjunto cobre os oito cantos da caixa $(\sigma,\omega,\eta_I)$ disciplinada. A coluna de fadiga s\'o acima de 40h fica abaixo da penalidade bilateral porque n\~ao penaliza a efici\^encia abaixo de 40 horas.}
\end{minipage}
\end{table*}

\begin{figure*}[t]
\centering
\includegraphics[width=0.92\textwidth]{fig_areq_vs_hours_pt.pdf}
\caption{Ganho requerido de PTF $A_{\text{req}}$ como fun\c{c}\~ao do teto semanal de horas.}\label{fig:areq_curve}
\vspace{2pt}
\begin{minipage}{0.92\textwidth}
\justifying
{\footnotesize \textit{Notas:} $A_{\text{req}}$ \'e o ganho \'unico de PTF que mant\'em o produto agregado constante no teto indicado. Os marcadores verticais distinguem a alternativa de 40h na proposta do relator do endpoint de 36h da PEC 8/2025 original. A curva de fadiga s\'o acima de 40h fornece $6{,}63\%$ no teto de 36h; a curva de penalidade bilateral fornece $8{,}18\%$. As duas curvas coincidem acima de $h^{*}=40$ e divergem abaixo do pico, refletindo a hip\'otese distinta sobre $e(h)$ abaixo do pico de efici\^encia.}
\end{minipage}
\end{figure*}

A figura tamb\'em ajuda a ler o que o modelo est\'a dizendo. As primeiras quatro horas de redu\c{c}\~ao s\~ao baratas em unidades de PTF porque o canal de fadiga trabalha a favor da reforma. As quatro horas seguintes s\~ao caras porque a reforma passa a retirar trabalho efetivo. O ap\^endice mostra o diagn\'ostico por tr\'as dessa afirma\c{c}\~ao: sem nenhum canal de efici\^encia das horas, o endpoint de 36 horas exigiria $7{,}74\%$ de PTF; deslocar o pico de efici\^encia em apenas uma hora muda mais o endpoint do que v\'arios par\^ametros secund\'arios. Por isso o m\'aximo de bem-estar perto de 41--42 horas deve ser lido como resultado da agenda calibrada de efici\^encia, n\~ao como regra de pol\'itica independente do modelo.

Ancoramos $A_{\text{req}}$ no hist\'orico brasileiro de PTF usando a s\'erie \texttt{rtfpna} da PWT 11.0 \citep{FeenstraInklaarTimmer2015,PWT2025,FREDPWTBrazil2026}. Essa compara\c{c}\~ao \'e uma r\'egua de escala, n\~ao uma valida\c{c}\~ao causal do contrafactual, pois o res\'iduo macro anual da PWT n\~ao \'e o mesmo objeto que a resposta de produtividade de curto prazo do modelo. A PTF brasileira caiu no per\'iodo p\'os-1990, ent\~ao um ganho \'unico de seis a oito por cento \'e grande nessa r\'egua. A alternativa de 40 horas est\'a em outra regi\~ao da escala. O ap\^endice traduz os ganhos \'unicos em taxas anualizadas para horizontes de faseamento de tr\^es, cinco e dez anos; essa leitura de transi\c{c}\~ao \'e mais informativa do que perguntar se uma d\'ecada hist\'orica espec\'ifica ``repete'' o contrafactual.

\begin{figure*}[t]
\centering
\includegraphics[width=0.95\textwidth]{fig_areq_vs_tfp_history_pt.pdf}
\caption{PTF brasileira no longo prazo.}\label{fig:tfp_history}
\vspace{2pt}
\begin{minipage}{0.95\textwidth}
\justifying
{\footnotesize \textit{Notas:} A figura plota a s\'erie \texttt{rtfpna} da PWT 11.0 (PTF real a pre\c{c}os nacionais constantes) para o Brasil. A linha pontilhada marca o n\'ivel de 1990; a faixa sombreada marca a melhor d\'ecada m\'ovel p\'os-1990 na s\'erie atualizada. \textit{Fonte:} Penn World Table 11.0 via FRED, s\'erie \texttt{RTFPNABRA632NRUG} \citep{FeenstraInklaarTimmer2015,PWT2025,FREDPWTBrazil2026}.}
\end{minipage}
\end{figure*}

Um teste de estresse pessimista no ap\^endice pergunta o que ocorre se a reorganiza\c{c}\~ao reduz temporariamente a produtividade. Esse caso n\~ao \'e a previs\~ao central; ele serve como guarda-corpo. Uma disrup\c{c}\~ao de PTF de $1\%$ empurra o requisito do endpoint com penalidade bilateral para acima de nove por cento, mostrando que a aritm\'etica da reforma depende da qualidade de implementa\c{c}\~ao, e n\~ao apenas da oferta de horas no longo prazo.

% ==================================================================
% 5. Bem-estar e heterogeneidade
% ==================================================================
\subsection{Bem-estar e incid\^encia}\label{sec:welfare}

O c\'alculo de bem-estar deve ser lido como diagn\'ostico de consumo e lazer. Ele pergunta algo estreito: se emprego \'e fixo, sal\'arios n\~ao barganham e a reforma entrega mecanicamente mais tempo fora do trabalho aos incumbentes, o ganho de lazer compensa a perda de consumo? Isso n\~ao \'e um teorema de bem-estar social. Sa\'ude, produ\c{c}\~ao dom\'estica, barganha, pre\c{c}os, transi\c{c}\~oes de emprego e pesos distributivos ficam fora do c\'alculo.

Dentro desse objeto estreito, a curva tem um formato claro. Perto de 41 a 42 horas, o indicador compra lazer a baixo custo de produto. Abaixo de 40 horas, o pre\c{c}o do lazer adicional sobe porque a reforma j\'a n\~ao est\'a apenas cortando fadiga de jornadas longas. Por isso o indicador de agente representativo fica negativo antes do endpoint. As especifica\c{c}\~oes de fadiga s\'o acima de 40h e penalidade bilateral concordam na regi\~ao de redu\c{c}\~ao moderada e divergem principalmente abaixo do pico de efici\^encia.

\begin{figure*}[t]
\centering
\includegraphics[width=0.92\textwidth]{fig_welfare_schedule_pt.pdf}
\caption{Curva de bem-estar por teto de horas sob ambas as calibra\c{c}\~oes.}\label{fig:welfare_schedule}
\vspace{2pt}
\begin{minipage}{0.92\textwidth}
\justifying
{\footnotesize \textit{Notas:} Cada curva plota o bem-estar agregado $\Delta$CV da Equa\c{c}\~ao~\eqref{eq:cv} em tetos sucessivamente menores, em rela\c{c}\~ao ao baseline de 44h, sob as especifica\c{c}\~oes de fadiga s\'o acima de 40h e penalidade bilateral. O c\'alculo usa prefer\^encias GHH de agente representativo; n\~ao \'e uma ordena\c{c}\~ao de bem-estar social. $\Delta$CV atinge m\'aximo perto de 42h em cerca de $+0{,}62\%$ e torna-se negativo abaixo de 40h. No endpoint de 36h, o bem-estar cai $1{,}76\%$ sob fadiga s\'o acima de 40h e $3{,}63\%$ sob penalidade bilateral.}
\end{minipage}
\end{figure*}

Um ganho de PTF de aproximadamente $1{,}4\%$ no teto de 36h \'e suficiente para tornar o endpoint neutro no indicador GHH sob fadiga s\'o acima de 40h. O limiar sob penalidade bilateral \'e $2{,}9\%$. Ambos os limiares ficam bem abaixo do $A_{\text{req}}$ neutro em produto porque o lazer compensa parcialmente a perda de consumo dentro desse objeto. O ap\^endice mostra que o sinal diagn\'ostico sobrevive a prefer\^encias separ\'aveis log e CRRA(2): a alternativa de 40h permanece positiva, enquanto o endpoint de 36h permanece negativo.

O n\'umero agregado de bem-estar mascara incid\^encia. Na decomposi\c{c}\~ao por tipo sob penalidade bilateral reportada no ap\^endice online, trabalhadores formais que permanecem formais ganham $+1{,}41\%$ em varia\c{c}\~ao compensat\'oria de primeira ordem porque trabalham cerca de seis horas a menos. A margem de trabalhadores empurrada para a informalidade perde $-28{,}57\%$. Esse grupo \'e pequeno, $1{,}9\%$ dos trabalhadores, mas sua perda \'e grande. Trabalhadores informais n\~ao afetados pelo teto estatut\'ario t\^em incid\^encia direta zero no modelo. O valor de agente representativo da Equa\c{c}\~ao~\eqref{eq:cv} difere da soma por tipo ponderada por participa\c{c}\~oes porque o composto GHH \'e n\~ao linear. Por isso, o m\'aximo de bem-estar em 41 a 42h deve ser lido como resultado interno ao modelo de agente representativo, com incid\^encia distributiva reportada separadamente.

O schedule do indicador GHH n\~ao implica oposi\c{c}\~ao a redu\c{c}\~oes da jornada. Ele mostra que, sob a calibra\c{c}\~ao representativa do modelo, a margem que eleva esse indicador est\'a entre 41 e 42 horas, n\~ao em 36. O schedule tamb\'em separa a perda de produto sob produtividade fixa, o ganho de produtividade que compensaria essa perda e o diagn\'ostico de consumo-lazer. Esses objetos n\~ao se movem juntos. A regi\~ao em que o indicador sobe \'e mais estreita do que a regi\~ao neutra em produto.

Custos de ajuste e informalidade diferem entre tamanhos de firma e setores, produzindo incid\^encia heterog\^enea da reforma. Firmas pequenas ($\leq 49$ empregados) requerem um ganho de PTF maior do que firmas grandes, refletindo cunhas de formaliza\c{c}\~ao maiores e informalidade de base mais alta. Aplicar o arcabou\c{c}o nacional separadamente a tr\^es setores amplos com par\^ametros setoriais da PNAD entrega um $A_{\text{req}}$ por setor de $1{,}3$ a $1{,}8\%$ para a alternativa de 40 horas e $6{,}91$ a $8{,}05\%$ para o endpoint de 36 horas sob penalidade bilateral. A informalidade sobe em todos os setores no endpoint de 36 horas. Esses s\~ao contrafactuais de equil\'ibrio parcial setor a setor, n\~ao uma an\'alise de incid\^encia insumo-produto. O $A_{\text{req}}$ da ind\'ustria supera ligeiramente o da agricultura porque o moderado $\eta_I=0{,}40$ combinado \`a base formal maior da ind\'ustria amortece o gap mec\^anico entre os dois setores. Os servi\c{c}os concentram a maior parte da perda agregada de produto via sua escala de produto, n\~ao por efeito maior por trabalhador. O ap\^endice online reporta a decomposi\c{c}\~ao completa e as sensibilidades de par\^ametros.

% ==================================================================
% 5. Implica\c{c}\~oes de pol\'itica p\'ublica
% ==================================================================
\section{Implica\c{c}\~oes para o desenho da reforma}\label{sec:conclusion}

O objeto relevante de pol\'itica j\'a n\~ao \'e apenas a vers\~ao original de 36 horas da PEC 8/2025. O debate vivo tamb\'em inclui alternativas de 40 horas, escala 5x2 e transi\c{c}\~oes faseadas. Por isso, o modelo calcula a sequ\^encia completa de tetos de horas, e n\~ao apenas o endpoint mais intenso.

Um teto de 40 horas \'e uma reforma moderada nessa aritm\'etica. Ele exige cerca de $2\%$ de PTF para restaurar o produto e deixa positivo o indicador GHH representativo. O movimento direto de 44 para 36 horas \'e maior. Ele afeta a grande maioria dos trabalhadores formais brasileiros e reduz mecanicamente o trabalho efetivo da economia em cerca de $14\%$. No endpoint de 36 horas, a perda de produto \'e compensada por um salto \'unico de produtividade entre $6{,}63\%$ sob fadiga s\'o acima de 40h e $8{,}18\%$ sob penalidade bilateral.

Esses n\'umeros s\~ao grandes quando lidos como benchmark de escala contra o hist\'orico brasileiro p\'os-1990 atualizado. Na PWT 11.0, a PTF brasileira cai de 1990 a 2023 e tamb\'em no horizonte pr\'e-Covid de 1990 a 2019. Mesmo a melhor d\'ecada m\'ovel p\'os-1990, de 2000 a 2010, cresce apenas $+0{,}05\%$ ao ano. Em um faseamento de cinco anos, o endpoint de 36h sob fadiga s\'o acima de 40h exigiria $1{,}29\%$ ao ano de PTF, enquanto a alternativa de 40h exigiria $0{,}38\%$ ao ano. Essa compara\c{c}\~ao n\~ao prova inviabilidade; ela apenas coloca o tamanho do ganho requerido na escala da experi\^encia recente brasileira.

O schedule \'e inclinado. Mover de 44 para 40 horas fica na ordem de metas moderadas de m\'edio prazo. Mover diretamente para 36 horas imp\~oe uma exig\^encia elevada de produtividade pelo padr\~ao brasileiro p\'os-1990. Sem esse salto, o modelo prev\^e perda de produto de curto prazo de $6{,}6$ a $8{,}0$ por cento e queda do indicador agregado GHH de $1{,}8$ a $3{,}6$ por cento no endpoint de 36 horas. O requisito de produtividade \'e um benchmark, n\~ao uma previs\~ao.

Uma redu\c{c}\~ao moderada do teto eleva o indicador GHH dentro do modelo. Reduzir a jornada de 44 para 41 a 42 horas aumenta esse diagn\'ostico porque o valor atribuido ao lazer adicional supera o consumo perdido na calibra\c{c}\~ao representativa. A alternativa de 40 horas fica perto da borda da regi\~ao positiva. Abaixo de 40h, o indicador fica negativo. O teto que maximiza esse diagn\'ostico est\'a perto de 42 horas, com 41h essencialmente indistingu\'ivel. Esse m\'aximo interior \'e uma propriedade do modelo, n\~ao uma recomenda\c{c}\~ao de pol\'itica.

O custo do endpoint \'e concentrado. Firmas pequenas, que empregam a maior parte dos trabalhadores formais no Brasil, carregam o maior requisito de produtividade por trabalhador. Ind\'ustria e agricultura carregam os maiores requisitos por setor. Todos os tr\^es setores veem a informalidade subir. Servi\c{c}os concentra a maior parte da perda agregada de produto devido \`a sua participa\c{c}\~ao no emprego e no produto. O exerc\'icio setorial \'e de equil\'ibrio parcial, setor por setor, e n\~ao uma an\'alise insumo-produto de incid\^encia.

A informalidade sobe, mas n\~ao explica a maior parte do custo. O modelo calibrado, disciplinado pelo pr\^emio salarial formal-informal observado, encontra um aumento agregado de informalidade de $1{,}6$ a $1{,}9$ p.p. Essa eleva\c{c}\~ao \'e material para a contabilidade de bem-estar e pequena em rela\c{c}\~ao \`a contra\c{c}\~ao mec\^anica de aproximadamente $15\%$ das horas. O custo de produto vem sobretudo de menos horas trabalhadas. Ainda assim, a reforma empurra trabalhadores adicionais para a informalidade, enfraquecendo os objetivos formalizadores da pol\'itica.

A Figura~\ref{fig:transition_map} retoma o mapa de transi\c{c}\~ao da vers\~ao anterior do artigo com a calibra\c{c}\~ao atual, mas foca a alternativa de 44 para 40 horas. Ela cruza a facilidade de migra\c{c}\~ao formal-informal, medida por $\sigma_{\text{sub}}$, com um al\'ivio tempor\'ario no wedge das pequenas firmas nesse cen\'ario de teto moderado. O exerc\'icio n\~ao \'e uma recomenda\c{c}\~ao de subs\'idio espec\'ifico; ele mostra que instrumentos de transi\c{c}\~ao pr\'o-pequenas firmas podem reduzir simultaneamente o vazamento para a informalidade e o ganho de PTF requerido.

\begin{figure*}[t]
\centering
\includegraphics[width=0.88\textwidth]{fig_transition_map_pt.pdf}
\caption{Mapa de transi\c{c}\~ao do $A_{\text{req}}$ para 44--40h: facilidade de migra\c{c}\~ao e al\'ivio no wedge nas pequenas firmas.}\label{fig:transition_map}
\vspace{2pt}
\begin{minipage}{0.88\textwidth}
\justifying
{\footnotesize \textit{Notas:} O eixo vertical mede o al\'ivio no wedge $1-\tau_S^{cap}/\tau_S$: quanto maior o valor, maior a redu\c{c}\~ao do custo de formaliza\c{c}\~ao nas pequenas firmas no cen\'ario de redu\c{c}\~ao de 44 para 40 horas (por exemplo, $0{,}5=$ corte de $50\%$). Cores mais claras indicam maior $A_{\text{req}}$; as linhas brancas s\~ao iso-$A_{\text{req}}$ (valores em \%). A figura reimplementa o mapa da vers\~ao anterior usando o c\'odigo e a calibra\c{c}\~ao atuais do pacote de replica\c{c}\~ao.}
\end{minipage}
\end{figure*}

A interpreta\c{c}\~ao exige separar produto, PTF requerida e bem-estar. Uma redu\c{c}\~ao moderada pode reduzir o produto, elevar o bem-estar e ainda exigir um ganho de PTF positivo, embora pequeno. Esse \'e o caso entre 44 e 42 horas no modelo.

\paragraph{Por que um pacote de transi\c{c}\~ao importa?} No exerc\'icio de 44 para 40 horas, o pacote n\~ao \'e um resgate macro para viabilizar um salto grande de produtividade; ele \'e um seguro de implementa\c{c}\~ao. No ponto central do mapa, o ganho requerido de PTF cai de cerca de $2{,}0\%$ sem al\'ivio no wedge das pequenas para aproximadamente $1{,}0\%$ com al\'ivio de $20\%$, $0{,}6\%$ com al\'ivio de $30\%$ e pr\'oximo de zero com al\'ivio de $50\%$. A interpreta\c{c}\~ao substantiva \'e que instrumentos tempor\'arios de transi\c{c}\~ao podem transformar uma meta moderada de produtividade em uma meta pequena, justamente na margem em que o modelo prev\^e maior fragilidade: pequenas firmas com maior custo de permanecer formais. Evid\^encia recente de pilotos de semana de quatro dias refor\c{c}a essa leitura: \citet{FanSchorKellyGu2025} analisam uma interven\c{c}\~ao de seis meses em 141 organiza\c{c}\~oes e 2.896 trabalhadores em seis pa\'ises, com redu\c{c}\~ao de jornada sem corte salarial, e encontram melhora em burnout, satisfa\c{c}\~ao, sa\'ude mental e sa\'ude f\'isica. Esse resultado deve ser lido como evid\^encia de viabilidade organizacional e bem-estar; a medida de produtividade ali \'e sobretudo capacidade de trabalho auto-reportada, n\~ao PTF agregada, de modo que complementa, mas n\~ao substitui, a aritm\'etica de $A_{\mathrm{req}}$.

O modelo \'e de curto prazo e equil\'ibrio parcial. O ajuste end\'ogeno do estoque de capital em horizonte de tr\^es a cinco anos reduz o ganho de produtividade requerido em cerca de um ponto e meio percentual. O emprego total \'e mantido fixo, de modo que o modelo desliga a partilha de trabalho. A evid\^encia emp\'irica sobre partilha de trabalho na Alemanha \citep{Hunt1996} e na reforma constitucional brasileira de 1988 \citep{GonzagaMenezesFilhoCamargo2003} sugere que essa restri\c{c}\~ao \'e compat\'ivel com a evid\^encia dispon\'ivel, mas n\~ao inocente. A experi\^encia portuguesa de 1996 \citep{AsaiLopesTondini2024} documenta ganhos de produtividade a n\'ivel de firma maiores do que os que o modelo produz, o que sugere que o modelo faz uma leitura conservadora.

Sob esses limites, a aritm\'etica da produtividade torna um caminho faseado mais f\'acil de absorver do que um salto direto para 36 horas. Uma redu\c{c}\~ao com parada intermedi\'aria pr\'oxima de 40 a 42 horas \'e consistente com o diagn\'ostico GHH do modelo, mant\'em o requisito de produtividade muito abaixo do endpoint de 36 horas e ganha tempo para o estoque de capital, a margem informal e a produtividade a n\'ivel de firma se ajustarem antes do pr\'oximo passo. A implica\c{c}\~ao \'e condicional ao modelo: tamanho e ritmo da redu\c{c}\~ao importam para a aritm\'etica de produto e para o indicador de consumo-lazer representativo.

\FloatBarrier

% === Ap\^endice online, c\'odigo e dados / Bibliografia ===
\paragraph{Ap\^endice online.} Um ap\^endice online com fontes de dados, omiss\~oes do equil\'ibrio parcial, detalhes da calibra\c{c}\~ao, grades de sensibilidade, decomposi\c{c}\~ao setorial, bem-estar por tipo, decomposi\c{c}\~ao do produto, tabela hist\'orica de PTF, deriva\c{c}\~oes e robustez da efici\^encia das horas est\'a dispon\'ivel no \href{\onlineappendixurl}{reposit\'orio p\'ublico de replica\c{c}\~ao}.

\paragraph{Disponibilidade de c\'odigo e dados.} C\'odigo, targets de calibra\c{c}\~ao e testes unit\'arios para reproduzir cada tabela e figura numerada neste artigo est\~ao dispon\'iveis no \href{https://github.com/vr-rodrigues/jornada-6x1-replication}{reposit\'orio p\'ublico do autor}.

\bibliographystyle{aer}
\bibliography{bibliography_clean}

\clearpage
\onecolumn
\section*{Informa\c{c}\~oes de submiss\~ao}
\setlength{\parindent}{0pt}
\setlength{\parskip}{4pt}

\textbf{T\'itulo em ingl\^es (English title)}

Public Policy Note: How much productivity do we need to reduce working hours?

\section*{Abstract}

Brazil's working-time debate is no longer only a choice between keeping the 44-hour week and moving directly to 36 hours. Alternatives around 40 hours, a five-day schedule and phased transitions are also on the table. This policy note asks a simple question for that choice: how much more productive would the economy need to become for each option not to reduce output in the short run? To answer, I combine Brazilian data on hours worked, informality, firm size and sectoral composition with a model of adjustment between formal and informal employment. The main result is that a move to 40 hours requires a productivity gain of about 2 percent. A direct move to 36 hours requires a much larger jump, between 6.6 and 8.2 percent, which is high relative to Brazil's recent productivity record. Informality also rises in the 36-hour scenario, by about 1.6 to 1.9 percentage points, but the main cost comes from fewer formal hours worked. The exercise does not say whether the reform should or should not move forward; it shows that size, timing and transition instruments change the arithmetic substantially. For policymakers, the message is direct: a phased route, with a stop near 40 hours, requires a much smaller productivity target than an immediate jump to 36 hours.

\textbf{Palavras-chave em ingl\^es (Keywords)}

Working hours; informality; productivity; Brazil; policy design; transition.

\textbf{ORCID}

Victor Rangel: \href{https://orcid.org/0000-0002-4520-2795}{https://orcid.org/0000-0002-4520-2795}.

\textbf{Declara\c{c}\~ao de contribui\c{c}\~ao dos autores}

Este manuscrito \'e de autoria \'unica; portanto, a declara\c{c}\~ao de contribui\c{c}\~ao dos autores n\~ao se aplica.

\textbf{Conflito de interesses (Conflict of interest)}

O autor declara n\~ao haver conflito de interesses financeiro, institucional ou pessoal relacionado a este manuscrito.

\textbf{Declara\c{c}\~ao de disponibilidade de dados de pesquisa (Data availability statement)}

O c\'odigo, os targets de calibra\c{c}\~ao, os testes unit\'arios, as figuras, as tabelas e as instru\c{c}\~oes de reprodu\c{c}\~ao est\~ao dispon\'iveis em \url{https://github.com/vr-rodrigues/jornada-6x1-replication}. Os dados brutos utilizados no artigo v\^em de bases p\'ublicas, incluindo PNAD Cont\'inua/IBGE, DIEESE, Penn World Table e fontes oficiais citadas no manuscrito, e podem ser regenerados pelos scripts do pacote de replica\c{c}\~ao.

\textbf{Dados e uso de IA}

Este manuscrito utilizou IA agentic, via Codex, como apoio t\'ecnico e editorial durante a prepara\c{c}\~ao do material de pesquisa. O autor definiu a pergunta, a estrat\'egia de modelagem, a calibra\c{c}\~ao, a interpreta\c{c}\~ao dos resultados e a reda\c{c}\~ao final. A autoria intelectual do projeto e a responsabilidade por eventuais erros s\~ao integralmente do autor.

\end{document}
